% Options for packages loaded elsewhere
\PassOptionsToPackage{unicode}{hyperref}
\PassOptionsToPackage{hyphens}{url}
\PassOptionsToPackage{dvipsnames,svgnames,x11names}{xcolor}
%
\documentclass[
  10pt,
  a4paper,
]{article}
\usepackage{amsmath,amssymb}
\usepackage{iftex}
\ifPDFTeX
  \usepackage[T1]{fontenc}
  \usepackage[utf8]{inputenc}
  \usepackage{textcomp} % provide euro and other symbols
\else % if luatex or xetex
  \usepackage{unicode-math} % this also loads fontspec
  \defaultfontfeatures{Scale=MatchLowercase}
  \defaultfontfeatures[\rmfamily]{Ligatures=TeX,Scale=1}
\fi
\usepackage[]{mathptmx}
\ifPDFTeX\else
  % xetex/luatex font selection
\fi
% Use upquote if available, for straight quotes in verbatim environments
\IfFileExists{upquote.sty}{\usepackage{upquote}}{}
\IfFileExists{microtype.sty}{% use microtype if available
  \usepackage[]{microtype}
  \UseMicrotypeSet[protrusion]{basicmath} % disable protrusion for tt fonts
}{}
\makeatletter
\@ifundefined{KOMAClassName}{% if non-KOMA class
  \IfFileExists{parskip.sty}{%
    \usepackage{parskip}
  }{% else
    \setlength{\parindent}{0pt}
    \setlength{\parskip}{6pt plus 2pt minus 1pt}}
}{% if KOMA class
  \KOMAoptions{parskip=half}}
\makeatother
\usepackage{xcolor}
\usepackage[margin=2.4cm]{geometry}
\usepackage{longtable,booktabs,array}
\usepackage{calc} % for calculating minipage widths
% Correct order of tables after \paragraph or \subparagraph
\usepackage{etoolbox}
\makeatletter
\patchcmd\longtable{\par}{\if@noskipsec\mbox{}\fi\par}{}{}
\makeatother
% Allow footnotes in longtable head/foot
\IfFileExists{footnotehyper.sty}{\usepackage{footnotehyper}}{\usepackage{footnote}}
\makesavenoteenv{longtable}
\setlength{\emergencystretch}{3em} % prevent overfull lines
\providecommand{\tightlist}{%
  \setlength{\itemsep}{0pt}\setlength{\parskip}{0pt}}
\setcounter{secnumdepth}{-\maxdimen} % remove section numbering
% pdflatex needs explicit declarations for the non-ASCII characters used here.
\usepackage[T1]{fontenc}
\usepackage[utf8]{inputenc}
\usepackage{textcomp}
\usepackage{amsmath,amssymb}
\DeclareUnicodeCharacter{2192}{$\rightarrow$}
\DeclareUnicodeCharacter{21D2}{$\Rightarrow$}
\DeclareUnicodeCharacter{2248}{$\approx$}
\DeclareUnicodeCharacter{00D7}{$\times$}
\DeclareUnicodeCharacter{03A9}{$\Omega$}
\DeclareUnicodeCharacter{2016}{$\|$}
\DeclareUnicodeCharacter{00B1}{$\pm$}
\DeclareUnicodeCharacter{00B5}{\textmu}
\DeclareUnicodeCharacter{00B7}{\textperiodcentered}
\DeclareUnicodeCharacter{00B2}{\textsuperscript{2}}
\DeclareUnicodeCharacter{00B0}{\textdegree}
\usepackage{etoolbox}
\usepackage{microtype}
\usepackage{booktabs}
\usepackage{longtable}
\usepackage{array}
\usepackage{ragged2e}
\usepackage{amsmath}
\usepackage{amssymb}
\usepackage[svgnames]{xcolor}

% Tables in this report are wide and text-heavy: shrink and left-align.
\AtBeginEnvironment{longtable}{\footnotesize\RaggedRight}
\setlength{\LTleft}{0pt}
\setlength{\LTright}{\fill}

% Block quotes as a marked-off recommendation box.
\let\oldquote\quote
\let\endoldquote\endquote
\renewenvironment{quote}
  {\begin{oldquote}\itshape\color{DarkSlateGray}}
  {\end{oldquote}}


\usepackage{fancyhdr}
\pagestyle{fancy}
\fancyhf{}
\fancyhead[L]{\footnotesize Analog Pad Buffers for IHP SG13CMOS5L}
\fancyhead[R]{\footnotesize\thepage}
\renewcommand{\headrulewidth}{0.4pt}

\setcounter{tocdepth}{2}
\setcounter{secnumdepth}{0}
\ifLuaTeX
  \usepackage{selnolig}  % disable illegal ligatures
\fi
\IfFileExists{bookmark.sty}{\usepackage{bookmark}}{\usepackage{hyperref}}
\IfFileExists{xurl.sty}{\usepackage{xurl}}{} % add URL line breaks if available
\urlstyle{same}
\hypersetup{
  colorlinks=true,
  linkcolor={DarkSlateGray},
  filecolor={Maroon},
  citecolor={Blue},
  urlcolor={MidnightBlue},
  pdfcreator={LaTeX via pandoc}}

\title{Analog Pad Buffers for IHP SG13CMOS5L}
\usepackage{etoolbox}
\makeatletter
\providecommand{\subtitle}[1]{% add subtitle to \maketitle
  \apptocmd{\@title}{\par {\large #1 \par}}{}{}
}
\makeatother
\subtitle{Design notes: figures of merit, reference topologies, and a bandwidth budget}
\author{Christoph Maier}
\date{3 September 2026}

\begin{document}
\maketitle

{
\hypersetup{linkcolor=DarkSlateGray}
\setcounter{tocdepth}{2}
\tableofcontents
}
\hypertarget{scope}{%
\subsection{1. Scope}\label{scope}}

Working notes toward a differential-input, analog-output pad driver (and
its dual, a single-ended-pin to differential on-chip input buffer) for
the IHP SG13CMOS5L open-source variant, in the context of a Chipalooza
submission. These notes cover:

\begin{itemize}
\tightlist
\item
  which figures of merit are invariant under bias-current scaling, and
  are therefore the ones worth arguing about;
\item
  what SG13CMOS5L removes relative to SG13G2, and how each omission
  constrains the topology choice;
\item
  a reference bibliography of class-AB output stages;
\item
  the relationship between the Monticelli and Hogervorst/Huijsing
  lineages;
\item
  a bandwidth budget with conservative / reasonable / aggressive tiers.
\end{itemize}

Nothing here is silicon-verified. Numbers marked \emph{estimate} are
order-of-magnitude reasoning, not simulation results.

\begin{center}\rule{0.5\linewidth}{0.5pt}\end{center}

\hypertarget{figures-of-merit-irrespective-of-bias-current}{%
\subsection{2. Figures of merit, irrespective of bias
current}\label{figures-of-merit-irrespective-of-bias-current}}

If \(W\) and \(I\) are scaled together at fixed inversion level, then
\(g_m\), GBW and SR all scale \(\propto I\) while noise power scales
\(\propto 1/I\). The quantities that survive that scaling are exactly
the dimensionless ratios and the per-current groups.

\hypertarget{genuinely-intensive-no-current-anywhere}{%
\subsubsection{2.1 Genuinely intensive (no current
anywhere)}\label{genuinely-intensive-no-current-anywhere}}

\begin{longtable}[]{@{}
  >{\raggedright\arraybackslash}p{(\columnwidth - 4\tabcolsep) * \real{0.3333}}
  >{\raggedright\arraybackslash}p{(\columnwidth - 4\tabcolsep) * \real{0.3333}}
  >{\raggedright\arraybackslash}p{(\columnwidth - 4\tabcolsep) * \real{0.3333}}@{}}
\toprule\noalign{}
\begin{minipage}[b]{\linewidth}\raggedright
Quantity
\end{minipage} & \begin{minipage}[b]{\linewidth}\raggedright
Units
\end{minipage} & \begin{minipage}[b]{\linewidth}\raggedright
Comment
\end{minipage} \\
\midrule\noalign{}
\endhead
\bottomrule\noalign{}
\endlastfoot
\(g_m/I_D\) & V\(^{-1}\) & Efficiency coordinate; everything else is
downstream of where you sit on this curve \\
\((g_m/I_D)\cdot f_T\) & Hz/V & The honest speed--efficiency invariant;
peaks in moderate inversion, \(IC \approx 0.1\ldots1\) \\
\(g_m r_o\) & --- & Intrinsic gain; sets static gain error
\(1/(\beta A_0)\). In 130 nm: \(\approx 25\ldots30\) at
\(L_\mathrm{min}\), \(100\ldots300\) at \(L \approx 1\) µm \\
\(I_\mathrm{out,peak}/I_Q\) & --- & Class-AB current gain. Class A pins
this at 1. Decent AB stages: 10--100 \\
\(V_\mathrm{out,pp}/V_\mathrm{DD}\) & --- & Swing efficiency,
i.e.~\(\sum V_\mathrm{DSsat}\) lost \\
\(A_{VT} = V_\mathrm{os}\sqrt{WL}\) & mV·µm & Matching constant,
\(\approx 3\ldots5\) mV·µm thin-oxide; useful normalized form is
\(V_\mathrm{os}/V_\mathrm{out,pp}\) \\
CMRR, PSRR & dB & At the frequency of interest \\
GBW\(/f_T\) & --- & Technology utilization; usually humiliating, and
usefully so \\
\(\mathrm{SR}/(\mathrm{GBW}\cdot V_\mathrm{step})\) & --- & \(>2\pi\)
means settling is linear rather than slew-limited \\
\(\eta = P_\mathrm{load}/P_\mathrm{supply}\) & --- & 25 \% ceiling for
class A into a resistive load \\
\end{longtable}

\hypertarget{per-current-groups}{%
\subsubsection{2.2 Per-current groups}\label{per-current-groups}}

Invariant, but only comparable at equal \(C_L\) and equal swing:

\begin{itemize}
\tightlist
\item
  \(\mathrm{FoM_{SS}} = \mathrm{GBW}\cdot C_L / I_\mathrm{DD}\)
  {[}MHz·pF/mA{]}
\item
  \(\mathrm{FoM_{LS}} = \mathrm{SR}\cdot C_L / I_\mathrm{DD}\)
  {[}V·pF/(µs·mA){]}
\item
  NEF / PEF, or the raw invariant
  \(v_{n,\mathrm{in}}^2 \cdot I_\mathrm{DD}\)
\item
  HD3 or SFDR, quoted at fixed \(V_\mathrm{out,pp}/V_\mathrm{DD}\)
  \textbf{and} fixed \(g_m/I_D\) --- otherwise it measures the overdrive
  choice, not the topology
\end{itemize}

\hypertarget{pad-specific-and-where-such-buffers-actually-fail}{%
\subsubsection{2.3 Pad-specific, and where such buffers actually
fail}\label{pad-specific-and-where-such-buffers-actually-fail}}

\begin{itemize}
\tightlist
\item
  \textbf{Transparency}: \(C_\mathrm{in}/C_\mathrm{tapped\ node}\), plus
  kickback charge per unit input swing. For a monitor buffer on a
  low-current node this dominates everything else --- the buffer that
  measures the DAC must not be the DAC's dominant load.
\item
  \textbf{Stability across the load corner}: phase margin over \(C_L\)
  from \textasciitilde5 pF (bare pad + ESD) to \textasciitilde50 pF +
  coax. A single-number GBW is meaningless without the load-range
  envelope.
\item
  \(|Z_\mathrm{out}(f)|\) and back-drive tolerance.
\item
  Oxide field at the pad under ESD clamp turn-on.
\end{itemize}

\hypertarget{the-one-that-is-not-a-ratio}{%
\subsubsection{2.4 The one that is not a
ratio}\label{the-one-that-is-not-a-ratio}}

\textbf{Area.} Offset and matching drag it in through
\(A_{VT}/\sqrt{WL}\), so any FoM table that omits area is quietly
comparing a 40 µm² driver to a 4000 µm² one.

\begin{center}\rule{0.5\linewidth}{0.5pt}\end{center}

\hypertarget{what-sg13cmos5l-takes-away}{%
\subsection{3. What SG13CMOS5L takes
away}\label{what-sg13cmos5l-takes-away}}

SG13CMOS5L is not SG13G2 minus the HBTs. Per IHP's open-source page, it
carries the SG13CMOS front end but has \textbf{no isolated NMOS, no MIM
capacitor, four thin metals and a single 2 µm TopMetal}. Three of those
hit an output driver directly.

\hypertarget{no-mim-capacitor}{%
\subsubsection{3.1 No MIM capacitor}\label{no-mim-capacitor}}

The compensation capacitor becomes a thick-oxide MOScap in accumulation
or a MOM comb on M1--M4. With only four thin metals, MOM density is poor
and the capacitor sits low in the stack where it couples to substrate. A
MOScap is voltage-dependent even in accumulation.

\textbf{Consequence:} any topology whose stability depends on precise
pole--zero cancellation --- Miller with nulling resistor, nested Miller,
active-feedback compensation --- becomes a corner-analysis problem. This
is the strongest single steer in the whole exercise:

\begin{quote}
Prefer \textbf{load-compensated single-stage} architectures, where the
dominant pole is \(C_L\) at the pad and \(C_c\) is either absent or
non-critical. Folded cascode or current-mirror OTA with a class-AB
output --- not a two-stage Miller amplifier.
\end{quote}

It also devalues the multistage-compensation literature (Leung \& Mok
damping-factor control, Ho/Mok active feedback), which assumes a
well-characterized linear \(C_c\).

\hypertarget{no-deep-n-well-isolated-nmos}{%
\subsubsection{3.2 No deep n-well / isolated
NMOS}\label{no-deep-n-well-isolated-nmos}}

Every NMOS bulk is the global p-substrate, shared with whatever digital
neighbour the shuttle places next door.

\begin{itemize}
\tightlist
\item
  NMOS source followers take full body effect with no escape ⇒ PMOS
  devices in their own n-well become the preferred followers and level
  shifters.
\item
  Substrate-coupled noise enters the output stage's NMOS side unfiltered
  ⇒ keep the signal path PMOS-dominant where possible, and \textbf{carry
  the signal off-chip differentially rather than single-ended} if the
  measurement setup permits. If it must be single-ended at the pad,
  route a reference pad alongside as a pseudo-differential return and
  budget for the substrate noise that cannot be filtered.
\end{itemize}

\hypertarget{one-topmetal-2-uxb5m-no-tm2}{%
\subsubsection{3.3 One TopMetal, 2 µm, no
TM2}\label{one-topmetal-2-uxb5m-no-tm2}}

Supply distribution to the output stage and the ESD return path are both
thinner than in SG13G2. This caps peak drive current and makes IR drop
across the pad ring a real term in the swing budget. It also removes a
shielding layer: with four thin metals, one must be spent on a shield
plane over the input pair, which is expensive when M4 is also the best
MOM plate and the best local supply mesh.

\hypertarget{revised-fom-weighting}{%
\subsubsection{3.4 Revised FoM weighting}\label{revised-fom-weighting}}

Given the above, the discriminating figures for this process are:

\begin{enumerate}
\def\labelenumi{\arabic{enumi}.}
\tightlist
\item
  swing efficiency \(V_\mathrm{out,pp}/\sum V_\mathrm{DSsat}\) on the
  3.3 V thick-oxide rail;
\item
  phase margin across the full \(C_L\) envelope \textbf{with \(C_c\) at
  its ±30 \% MOScap corner};
\item
  PSRR and substrate-coupling rejection --- a first-class spec here, not
  an afterthought;
\item
  area, since MOM/MOScap compensation is area-hungry and shuttle slots
  are small.
\end{enumerate}

\(g_m/I_D\) and intrinsic gain still set the design point but no longer
differentiate between candidates; the compensation and isolation
constraints do.

\begin{center}\rule{0.5\linewidth}{0.5pt}\end{center}

\hypertarget{reference-designs-class-ab-output-stages}{%
\subsection{4. Reference designs --- class-AB output
stages}\label{reference-designs-class-ab-output-stages}}

\begin{longtable}[]{@{}
  >{\raggedright\arraybackslash}p{(\columnwidth - 6\tabcolsep) * \real{0.2500}}
  >{\raggedright\arraybackslash}p{(\columnwidth - 6\tabcolsep) * \real{0.2500}}
  >{\raggedright\arraybackslash}p{(\columnwidth - 6\tabcolsep) * \real{0.2500}}
  >{\raggedright\arraybackslash}p{(\columnwidth - 6\tabcolsep) * \real{0.2500}}@{}}
\toprule\noalign{}
\begin{minipage}[b]{\linewidth}\raggedright
\#
\end{minipage} & \begin{minipage}[b]{\linewidth}\raggedright
Design / paper
\end{minipage} & \begin{minipage}[b]{\linewidth}\raggedright
Venue
\end{minipage} & \begin{minipage}[b]{\linewidth}\raggedright
Link
\end{minipage} \\
\midrule\noalign{}
\endhead
\bottomrule\noalign{}
\endlastfoot
1 & Monticelli, \emph{A quad CMOS single-supply op amp with rail-to-rail
output swing} --- the floating-battery AB mesh & JSSC SC-21(6), 1986,
1026--1034 &
\href{http://class.ece.iastate.edu/djchen/ee501/2008/MonticelliRailToRailOutSwing.pdf}{PDF
(Iowa State)} ·
\href{https://doi.org/10.1109/jssc.1986.1052645}{doi:10.1109/JSSC.1986.1052645} \\
2 & Carvajal, Ramírez-Angulo et al., \emph{The flipped voltage follower}
--- FVF/DFVF taxonomy & TCAS-I 52, 2005, 1276--1291 &
\href{https://doi.org/10.1109/TCSI.2005.851387}{doi:10.1109/TCSI.2005.851387} \\
3 & Sawigun \& Demosthenous, \emph{A compact rail-to-rail class-AB CMOS
buffer with slew-rate enhancement} & TCAS-II, 2012 &
\href{https://ieeexplore.ieee.org/document/6236111/}{IEEE Xplore} \\
4 & \emph{High-speed rail-to-rail class-AB buffer with compact adaptive
biasing for FPD} & MDPI Electronics 9(12), 2020 &
\href{https://www.mdpi.com/2079-9292/9/12/2018}{open access} \\
5 & Ramírez-Angulo / Carvajal et al., \emph{Class-AB output stages with
accurate quiescent current control by dynamic biasing} & Analog Integr.
Circ. Sig. Process. &
\href{https://link.springer.com/article/10.1023/A:1024453731969}{Springer} \\
6 & \emph{Low-voltage class-AB output stages using floating capacitors}
--- the sub-\(2V_{GS}\) alternative to Monticelli & Electron. Lett.
lineage &
\href{https://www.academia.edu/22725840/Low_voltage_class_AB_output_stages_for_CMOS_op_amps_using_floating_capacitors}{PDF} \\
7 & Hogervorst et al., \emph{Rail-to-rail constant-\(g_m\) input stage
and class-AB output stage} & Analog Integr. Circ. Sig. Process. &
\href{https://link.springer.com/article/10.1007/BF01239246}{Springer} \\
8 & \emph{Rail-to-rail low-power high-slew-rate CMOS analogue buffer}
--- \textless32 fF input capacitance & --- &
\href{https://www.academia.edu/10744243/Rail_to_rail_low_power_high_slew_rate_CMOS_analogue_buffer}{Academia} \\
9 & \emph{Rail-to-rail high-speed class-AB CMOS buffer with enhanced
slew rate} --- 1 nF at 3.5 µA \(I_Q\) & --- &
\href{https://www.researchgate.net/publication/274173822_A_Rail-To-Rail_Hign_Speed_Class-AB_CMOS_Buffer_with_Low_Power_and_Enhanced_Slew_Rate}{ResearchGate} \\
10 & López-Martín et al., \emph{The Flipped Voltage Follower: Theory and
Applications} & Springer LNEE &
\href{https://link.springer.com/chapter/10.1007/978-3-642-36329-0_12}{chapter} \\
\end{longtable}

\textbf{Bibliography only (link not verified; all appear in the
reference lists of the above):}

\begin{itemize}
\tightlist
\item
  F. You, S. H. K. Embabi, E. Sánchez-Sinencio, ``Low-voltage class-AB
  buffers with quiescent current control,'' \emph{JSSC} 33(6), 1998,
  915--920.
\item
  M. de Langen, J. H. Huijsing, \emph{JSSC} 33(10), 1998, 1482--1496.
\item
  G. Giustolisi, G. Palmisano, ``1.2-V op-amp with a dynamically biased
  output stage,'' \emph{JSSC} 35(4), 2000, 632--636.
\item
  M. D. Pardoen, M. G. Degrauwe, ``A rail-to-rail CMOS input/output
  power amplifier,'' \emph{JSSC} 25(4), 1990, 501--504.
\item
  J. A. Fisher, R. Koch, ``A highly linear CMOS buffer amplifier,''
  \emph{JSSC} 22, 1987, 330--334.
\item
  R. Hogervorst, J. P. Tero, R. G. H. Eschauzier, J. H. Huijsing, ``A
  compact power-efficient 3 V CMOS rail-to-rail input/output operational
  amplifier for VLSI cell libraries,'' \emph{JSSC} 29(12), Dec.~1994,
  1505--1513.
\end{itemize}

\hypertarget{reading-order-for-this-process}{%
\subsubsection{4.1 Reading order for this
process}\label{reading-order-for-this-process}}

Items 1 and 6 should be read against each other: Monticelli's mesh costs
more than \(2V_{GS}\) of headroom but needs no capacitor, whereas the
floating-capacitor variants buy headroom back by spending exactly the
component SG13CMOS5L does not provide. Items 3 and 4 are the useful pair
on slew enhancement. Item 2 is the compactness reference whose offset
and PSRR must be characterized independently, since the papers largely
do not.

\begin{center}\rule{0.5\linewidth}{0.5pt}\end{center}

\hypertarget{monticelli-vs.-hogervorst-huijsing-different-axes}{%
\subsection{5. Monticelli vs.~Hogervorst \& Huijsing --- different
axes}\label{monticelli-vs.-hogervorst-huijsing-different-axes}}

They are not the same topology, and the confusion is worth dissolving
because the two names sit on \textbf{different axes}.

\textbf{Monticelli (1986)} is about a class-AB \emph{output} stage plus
quiescent-current control. Its input stage is a conventional single
pair: the abstract claims an output swing to either rail together with
an input common-mode range that includes ground --- ground-sensing,
\emph{not} rail-to-rail input.

\textbf{Hogervorst \& Huijsing}'s headline contribution is precisely
what Monticelli lacks: the \textbf{constant-\(g_m\) rail-to-rail input
stage}, two complementary pairs with tail-current steering so
\(g_{m,\mathrm{tot}}\) stays flat as the common mode sweeps. They
publish output stages too, but the usual pairing in the literature is
\emph{Hogervorst input + some floating class-AB output}, and that output
can perfectly well be Monticelli's mesh. \textbf{They compose rather
than compete.}

\hypertarget{where-they-do-overlap}{%
\subsubsection{5.1 Where they do overlap}\label{where-they-do-overlap}}

Both output stages belong to Huijsing's \textbf{feedforward} class-AB
control family: a floating voltage source (translinear loop) between the
two output gates. This is distinct from \textbf{feedback} control with
an explicit minimum-current selector --- the de Langen \& Huijsing
branch (\emph{JSSC} 1998), which senses the smaller of the two output
currents and regulates it, giving tighter \(I_Q\) over corners at the
cost of a loop that can ring.

So Monticelli and the Hogervorst compact output stage are cousins; the
Delft feedback-controlled stages are a different animal.

\hypertarget{two-practical-differences}{%
\subsubsection{5.2 Two practical
differences}\label{two-practical-differences}}

\textbf{Headroom.} Monticelli's mesh needs roughly
\(2V_{GS} + 2V_\mathrm{DSsat}\) of stack. That is exactly the criticism
levelled by the floating-capacitor school, whose stated advantage is
accurate quiescent control at supplies near one threshold, where
Monticelli's requires more than \(2V_{GS}\). On a 3.3 V thick-oxide rail
with \(V_T \approx 0.5\ldots0.7\) V this is affordable; on the 1.2 V
core it is not.

\textbf{Drive asymmetry.} The mesh is not symmetric in the small-signal
path. One output device's gate is driven directly by the cascode
current; the complementary device is driven \emph{through} the mesh,
which behaves as an additional cascode stage with a pole set by
\(1/g_m\) of the mesh device and the capacitance at that node ---
contributing a non-dominant pole and extra delay on that path only. This
shows up as unequal rising/falling settling, and it interacts badly with
a compensation capacitor of uncertain value. Treat that pole's location
as a corner variable and simulate it explicitly rather than trusting a
nominal phase margin.

\begin{center}\rule{0.5\linewidth}{0.5pt}\end{center}

\hypertarget{hogervorst-et-al.-1994-as-a-donor-design}{%
\subsection{6. Hogervorst et al.~1994 as a donor
design}\label{hogervorst-et-al.-1994-as-a-donor-design}}

\emph{R. Hogervorst, J. P. Tero, R. G. H. Eschauzier, J. H. Huijsing,
``A Compact Power-Efficient 3 V CMOS Rail-to-Rail Input/Output
Operational Amplifier for VLSI Cell Libraries,'' JSSC 29(12), Dec.~1994,
1505--1513.}

Verdict: usable as an \textbf{output-stage donor, not as a cell to
copy}.

\hypertarget{keep-the-output-stage-and-its-biasing-figs.-5-8-11}{%
\subsubsection{6.1 Keep --- the output stage and its biasing (Figs. 5,
8, 11)}\label{keep-the-output-stage-and-its-biasing-figs.-5-8-11}}

The floating class-AB control \(M_{19}/M_{20}\) shifted \emph{into} the
folded-cascode summing circuit means the AB control contributes neither
noise nor offset: the two bias sources \(I_{b6}/I_{b7}\) that would
otherwise sit in parallel with the cascodes \(M_{14}/M_{16}\) are
eliminated outright. The floating current source \(M_{27}\)--\(M_{28}\),
built with the \emph{same} translinear structure as the AB control so
that its supply dependence cancels the control's, yields an \(I_Q\) that
does not walk with \(V_\mathrm{DD}\). A real solution to a real problem,
at a cost of two transistors.

Table I of that paper was measured at exactly 3.3 V --- the same rail as
IOVDD on the 5L shuttle --- with two stacked \(V_{GS}\) in the output
stage, which on IHP thick-oxide devices (\(V_T \approx 0.5\ldots0.7\) V)
leaves room.

\hypertarget{discard-the-constant-g_m-rail-to-rail-input-apparatus-section-ii}{%
\subsubsection{\texorpdfstring{6.2 Discard --- the constant-\(g_m\)
rail-to-rail input apparatus (Section
II)}{6.2 Discard --- the constant-g\_m rail-to-rail input apparatus (Section II)}}\label{discard-the-constant-g_m-rail-to-rail-input-apparatus-section-ii}}

For a differential \emph{on-chip} input at fixed, known common mode,
everything in their Figs. 1--4 is pure cost: the complementary pairs,
the two current switches, the 3× mirrors \(M_6\)--\(M_7\) and
\(M_9\)--\(M_{10}\), plus \(M_{29}\)--\(M_{31}\) whose only job is to
keep the positive-feedback loop from arming below 2.9 V.

What it buys in defects, from their own Table I and Section VI:

\begin{itemize}
\tightlist
\item
  CMRR collapses to \textbf{43 dB} in the take-over ranges, against 70
  dB elsewhere;
\item
  offset shifts about \textbf{2 mV per take-over range};
\item
  slew rate changes by a \textbf{factor of two} depending where the
  common mode sits (2 → 4 V/µs Miller; 4 → 8 V/µs cascoded-Miller).
\end{itemize}

A single PMOS pair at fixed common mode gives constant \(g_m\), constant
SR, no take-over artifacts, better CMRR, a bulk-tied-to-source device in
its own n-well (which matters with no deep n-well available), and
roughly 40 \% of the area back.

\hypertarget{the-porting-obstacle-c_m1c_m2}{%
\subsubsection{\texorpdfstring{6.3 The porting obstacle ---
\(C_{M1}/C_{M2}\)}{6.3 The porting obstacle --- C\_\{M1\}/C\_\{M2\}}}\label{the-porting-obstacle-c_m1c_m2}}

They are Miller capacitors tied from \(V_o\) to the output-device gates,
so the \textbf{full output swing stands across them}. With MIM that is
free; in 5L it means MOM fingers on four thin metals, or a thick-oxide
MOScap that will traverse accumulation-to-depletion as the output swings
--- a compensation capacitor that changes value with the signal it is
compensating.

Two mitigations available from the paper itself:

\begin{enumerate}
\def\labelenumi{\arabic{enumi}.}
\tightlist
\item
  The plain-Miller version has \textbf{66° phase margin} versus
  \textbf{53°} for cascoded-Miller. Take the Fig. 13 variant and spend
  the margin on capacitor uncertainty rather than on bandwidth.
\item
  They note cascoded-Miller can peak at high output current and
  recommend plain Miller above \textasciitilde3 mA --- which points the
  same way for a driver.
\end{enumerate}

\hypertarget{two-numbers-that-will-not-survive-the-port}{%
\subsubsection{6.4 Two numbers that will not survive the
port}\label{two-numbers-that-will-not-survive-the-port}}

\begin{itemize}
\tightlist
\item
  \textbf{Gain.} Their 85 dB came from 1 µm devices. On 130 nm
  thick-oxide with the same cascoding, budget 60--70 dB
  (\emph{estimate}); their own pointer to Bult \& Geelen gain boosting
  on \(M_{14}/M_{16}\) becomes less optional.
\item
  \textbf{Load.} 10 kΩ ‖ 10 pF at 3 mA peak is a \emph{chip-level}
  buffer. A metre of coax plus a scope front end is 100 pF and change.
  Scaling the output devices for that raises \(C_{GS,\mathrm{out}}\),
  drags the output pole down, and forces \(C_M\) up --- straight back
  into the capacitor problem. Size that loop before committing.
\end{itemize}

Their figure of merit, \(F = B/P_\mathrm{sup}\) at \(C_L = 10\) pF, is
supply-referred rather than current-referred; convert to
\(\mathrm{GBW}\cdot C_L/I_\mathrm{DD}\) to compare against anything
modern.

\hypertarget{incidental}{%
\subsubsection{6.5 Incidental}\label{incidental}}

The Fig. 6 → 8 → 9 → 10 → 11 progression, in which each figure repairs
one named defect of its predecessor, is unusually clean design-rationale
writing. If the shuttle wants documentation rather than only a GDS, that
structure is worth reusing as a template.

\begin{center}\rule{0.5\linewidth}{0.5pt}\end{center}

\hypertarget{bandwidth-budget}{%
\subsection{7. Bandwidth budget}\label{bandwidth-budget}}

\hypertarget{the-load-sets-everything-on-the-output-side}{%
\subsubsection{7.1 The load sets everything on the output
side}\label{the-load-sets-everything-on-the-output-side}}

For a load-compensated single stage,

\[\mathrm{GBW} = \frac{g_m}{2\pi C_L}\]

\(C_L\) is not the pad. It is pad (a few hundred fF) + ESD
(\textasciitilde1 pF) + bond wire + PCB trace + whatever instrument is
attached. A 10× scope probe alone is 10--15 pF. \textbf{Budget
\(C_L \approx 15\) pF} as the honest bench number, and 50--100 pF if
anyone ever hangs un-terminated coax on it.

At \(C_L = 15\) pF:

\begin{longtable}[]{@{}lll@{}}
\toprule\noalign{}
Target GBW & Required \(g_m\) & \(I_D\) at \(g_m/I_D = 10\)
V\(^{-1}\) \\
\midrule\noalign{}
\endhead
\bottomrule\noalign{}
\endlastfoot
10 MHz & 0.94 mS & ≈ 95 µA \\
100 MHz & 9.4 mS & ≈ 1 mA \\
\end{longtable}

Current is not the wall. \textbf{The non-dominant poles are.} A
thick-oxide 3.3 V device at \(L \approx 1\) µm in moderate inversion has
\(f_T\) of perhaps 1--3 GHz (\emph{estimate}), and the mirror and
cascode poles will land in the low hundreds of MHz. Pushing GBW to 100
MHz puts the unity-gain crossing uncomfortably close to them --- while
the compensation capacitor's value is uncertain because there is no MIM.

\hypertarget{slew-usually-binds-first}{%
\subsubsection{7.2 Slew usually binds
first}\label{slew-usually-binds-first}}

Full-power bandwidth requires \(\mathrm{SR} = 2\pi f V_\mathrm{pk}\),
i.e.~peak current \(I = C_L \cdot \mathrm{SR}\). For 1 V peak at 10 MHz
into 15 pF that is ≈ 0.94 mA of peak drive --- comfortable for a
class-AB stage running 50 µA quiescent. At 100 MHz it is 9.4 mA, which
on a single 2 µm TopMetal supply route starts consuming the swing budget
in IR drop.

\textbf{Quote small-signal GBW and full-power bandwidth separately.}
They are not the same spec, and the gap between them is the entire point
of class AB.

\hypertarget{calibration-tiers-output-buffer-differential-in-single-ended-pin}{%
\subsubsection{7.3 Calibration tiers --- output buffer (differential in
→ single-ended
pin)}\label{calibration-tiers-output-buffer-differential-in-single-ended-pin}}

\begin{longtable}[]{@{}
  >{\raggedright\arraybackslash}p{(\columnwidth - 6\tabcolsep) * \real{0.2500}}
  >{\raggedright\arraybackslash}p{(\columnwidth - 6\tabcolsep) * \real{0.2500}}
  >{\raggedright\arraybackslash}p{(\columnwidth - 6\tabcolsep) * \real{0.2500}}
  >{\raggedright\arraybackslash}p{(\columnwidth - 6\tabcolsep) * \real{0.2500}}@{}}
\toprule\noalign{}
\begin{minipage}[b]{\linewidth}\raggedright
Tier
\end{minipage} & \begin{minipage}[b]{\linewidth}\raggedright
GBW
\end{minipage} & \begin{minipage}[b]{\linewidth}\raggedright
FPBW at 1 V\(_\mathrm{pk}\)
\end{minipage} & \begin{minipage}[b]{\linewidth}\raggedright
Risk
\end{minipage} \\
\midrule\noalign{}
\endhead
\bottomrule\noalign{}
\endlastfoot
Conservative & 1--5 MHz & \textasciitilde1--2 MHz & first silicon works,
no drama \\
\textbf{Reasonable} & \textbf{10--30 MHz} & \textbf{5--10 MHz} & sweet
spot; poles comfortably clear \\
Aggressive & 50--100 MHz & 20 MHz & pole management becomes the whole
design; MOScap \(C_c\) uncertainty acute \\
Don't & \textgreater{} 200 MHz & --- & stop calling it a buffer \\
\end{longtable}

Above roughly 100 MHz through a wire-bonded pad into an unknown board,
the problem stops being ``amplifier'' and becomes ``transmission line'':
an impedance-matched 50 Ω driver is wanted, the load turns resistive,
and power rises by an order of magnitude. Different project.

\hypertarget{the-input-buffer-is-an-easier-problem}{%
\subsubsection{7.4 The input buffer is an easier
problem}\label{the-input-buffer-is-an-easier-problem}}

Single-ended pin → differential on-chip drives only internal capacitance
--- call it 200--500 fF. 50 MHz needs under 100 µS of \(g_m\). The
limits here are the ESD structure's series RC, the input pair's own
matching and noise, and the fact that deliberate band-limiting is
usually \emph{wanted} for anti-aliasing and noise.

\textbf{50--100 MHz is comfortable; 20 MHz is conservative.} The
asymmetry between the two directions is real --- roughly a decade of
headroom.

\hypertarget{two-arguments-for-staying-low}{%
\subsubsection{7.5 Two arguments for staying
low}\label{two-arguments-for-staying-low}}

\begin{itemize}
\tightlist
\item
  Noise scales as \(\sqrt{\mathrm{BW}}\). If the path characterizes
  anything quiet, excess bandwidth is a direct loss of resolution.
\item
  Keeping GBW modest lets \(C_L\) be the dominant pole and lets the
  Miller capacitor be dropped entirely --- the cleanest escape from the
  no-MIM problem of §3.1.
\end{itemize}

\hypertarget{working-stake}{%
\subsubsection{7.6 Working stake}\label{working-stake}}

\textbf{20 MHz out, 50 MHz in.} With the caveat that the right answer
depends on what is being observed: DC bias characterization needs 1 MHz;
observing a settling transient needs whatever bandwidth resolves the
settling constant of interest.

\begin{center}\rule{0.5\linewidth}{0.5pt}\end{center}

\hypertarget{fpbw-definition}{%
\subsection{8. FPBW --- definition}\label{fpbw-definition}}

\textbf{Full-Power Bandwidth}: the highest frequency at which the
amplifier still delivers its \emph{full rated output swing} without the
waveform being distorted by slew-rate limiting.

The derivation is one line. For
\(V_\mathrm{out} = V_\mathrm{pk}\sin(2\pi f t)\), the maximum rate of
change occurs at the zero crossing:

\[\left.\frac{dV}{dt}\right|_\mathrm{max} = 2\pi f V_\mathrm{pk}\]

The amplifier can only supply \(dV/dt\) up to its slew rate. Setting
them equal:

\[\boxed{\ \mathrm{FPBW} = \frac{\mathrm{SR}}{2\pi V_\mathrm{pk}}\ }\]

Above that frequency the output stops being a sine and becomes a
triangle wave: the amplifier is charging \(C_L\) with all the current it
has and cannot turn around fast enough. This is a hard, ugly,
large-signal nonlinearity, not a gentle roll-off.

\hypertarget{why-it-is-separate-from-gbw}{%
\subsubsection{8.1 Why it is separate from
GBW}\label{why-it-is-separate-from-gbw}}

GBW is a \textbf{small-signal} spec --- how fast the amplifier responds
to a wiggle small enough that everything stays linear. FPBW is
\textbf{large-signal} --- how fast it responds when asked for the whole
swing.

A class-A stage with fixed tail current \(I\) has
\(\mathrm{SR} = I/C_L\), and its FPBW can be a factor of ten or more
below its GBW. It then looks fast on a Bode plot and turns a 1 V sine
into a triangle. Closing that gap is exactly what class AB is for: on a
large step the output devices deliver many times the quiescent current,
so SR is no longer pinned to \(I_Q\).

The ratio \(\mathrm{SR}/(\mathrm{GBW}\cdot V_\mathrm{step})\) from §2.1
is the same idea in dimensionless form. Above about \(2\pi\), the
amplifier settles by linear exponential decay and the FPBW spec does not
constrain; below it, slewing dominates the settling time and GBW is a
fiction.

\hypertarget{consistency-check}{%
\subsubsection{8.2 Consistency check}\label{consistency-check}}

1 V peak at 10 MHz needs
\(\mathrm{SR} = 2\pi\times10^7\times1 \approx 63\) V/µs; into 15 pF,
\(I = C\cdot\mathrm{SR} \approx 0.94\) mA --- the same figure as §7.1.
The peak-current and FPBW framings are one calculation wearing two hats.

\begin{center}\rule{0.5\linewidth}{0.5pt}\end{center}

\hypertarget{open-questions-before-committing}{%
\subsection{9. Open questions before
committing}\label{open-questions-before-committing}}

\begin{enumerate}
\def\labelenumi{\arabic{enumi}.}
\tightlist
\item
  Does \texttt{ihp-sg13cmos5l} ship an \textbf{analog pad cell} --- an
  unbuffered pad with a characterized ESD network and known parasitic
  capacitance --- or only digital I/O cells? If digital-only, the choice
  is between designing an ESD structure (probably not sign-off-able on a
  shared shuttle) and living with a digital pad's diode clamps in the
  analog path.
\item
  What does the Chipalooza floorplan allow for \textbf{pad allocation},
  and is an analog supply domain separate from IOVDD available at all?
\item
  Is the far end of the link a \textbf{differential receiver}? If so, do
  not collapse to single-ended on chip (§3.2).
\item
  What is the actual \(C_L\) envelope of the intended bench setup?
  Everything in §7 hangs on it.
\end{enumerate}

\begin{center}\rule{0.5\linewidth}{0.5pt}\end{center}

\hypertarget{provenance-and-caveats}{%
\subsection{10. Provenance and caveats}\label{provenance-and-caveats}}

\begin{itemize}
\tightlist
\item
  The SG13CMOS5L device/BEOL inventory is taken from IHP's open-source
  request page, which described the open PDK as under development for
  early 2026. Verify against the \texttt{IHP-GmbH/ihp-sg13cmos5l}
  repository rather than the marketing page.
\item
  Items 1--10 in §4 have verified links. The bibliography-only list does
  not; those citations come from the reference lists of the linked
  papers.
\item
  The Monticelli-vs-Huijsing feedforward/feedback taxonomy in §5.1 is
  asserted from secondary sources. Huijsing's \emph{Operational
  Amplifier Theory and Design} carries the canonical version and should
  be checked before the distinction carries weight in a design
  justification.
\item
  Every number marked \emph{estimate} is order-of-magnitude reasoning.
  Nothing here has been simulated in SG13CMOS5L.
\end{itemize}

\end{document}
